\section{Literature Review}


\subsection{Party Brand}

Political parties function as brands that signal policy positions to average voters. As Lupu (2016, pp. 17--18) puts it: ``Voters develop perceptions of political parties through their observations of a party and its behavior \ldots{} they learn what to associate with the prototypical Democrat\ldots{} Those prototypes are what I refer to as party brands.'' Those brands provide informational shortcuts, allowing voters to infer party positions often without detailed policy knowledge. Voters form stronger attachments to parties whose brands are distinguishable, and convergence between rival parties weakens those attachments \citep{lupuPartyBrandsPartisanship2013}. Perceiving parties as polarized has the same effect in the opposite direction, since citizens become more partisan as they perceive polarization increasing \citep{lupuPartyPolarizationMass2015}. When brands dilute, partisanship weakens and parties can break down entirely \citep{lupuPartyBrandsCrisis2016}. To avoid brand dilution, parties are motivated to occupy distinct positions (between-party differentiation) and maintain consistent messaging (within-party coherence), as ``declining levels of ideological and issue conflict undercut the relevance of both parties and partisan ties `` \citep[p. 110]{schmittPoliticalPartiesDecline1998}.

Brands also evolve as parties strategically reposition ideologically and contest which issue organizes the current political conflict. This is often done through politicization and polarization with elite messaging. Politicization is a process by which parties problematize and reconstruct meanings of a particular issue to be more abstract and open for indetermination, resulting in contestations and an increase in voters' perceived saliency and importance \citep{mazzucaMakingItAbstract2023,slothuusWhenCanPolitical2010}. Take the Ukraine aid as an example. Parties start the process by setting the ``undetermined'' premise of the issue \citep{koselleckFuturesSemanticsHistorical2004}, such that whether to spend a limited amount of money on helping another country is something that should be openly discussed. Parties then frame the issue to be more abstract and open for interpretation \citep{borghiWordsSocialTools2014,Zurn2020}, such that aids for Ukraine could mean loss of domestic spending and national interest. This ``undetermined'' premise of the issue now creates new ideological (e.g., fund vs. no-fund) and saliency space for parties to compete and build up clearer branding \citep{eliasPositionSelectiveEmphasis2015,guinaudeauWhatIssueCompetition2014,slothuusMoreWeightingCognitive2008,Zurn2020}. Further competition on a politicized issue introduces issue ownership and elite polarization (i.e., polarization between political elites), and issue ownership leads to greater visibility and coverage \citep{campIssueOwnershipDeterminant2018}. Therefore, parties are incentivized to polarize to mobilize the limited voter capital.

For voters, party cues can reduce the information demands of evaluating complex policies. It is well-known now that voters often use ``cognitive heuristics'' to make various political decisions \citep{adamsCompanyYouKeep2016,albrightDoesPoliticalKnowledge2009,banerjeeInformedVotersMake2011,lauAdvantagesDisadvantagesCognitive2001}, from what information to look for \citep{danceyPartyIdentificationIssue2010} to which candidate to vote for \citep{downsEconomicTheoryPolitical1957}. In fact, being more politically sophisticated or informed does not guarantee better normative political decisions than simply relying on heuristics cues such as party brands, yet it can exponentially increase cognitive costs in a market that enables parties to profit from increasing competition and information \citep{banerjeeInformedVotersMake2011,lauAdvantagesDisadvantagesCognitive2001}. Nevertheless, the effect of party cues are constrained by competing identities, accuracy goals, ambivalence, contradictory information, and reality constraints \citep{leeperPoliticalPartiesMotivated2014}. Across nine United States foreign-policy experiments that held substantive information constant while varying its generic, Democratic, or Republican attribution, cue effects differed with baseline disagreement and partisan polarization across issues \citep{guisingerMappingBoundariesElite2017}. Elites therefore have reason to politicize and polarize policies to sharpen party brands, even though the direction in which voters respond to those brands varies across issues and contexts.

Furthermore, brands do not derive solely from direct party communications (e.g., speeches, platforms, and social media). They are also shaped by the perceived image created through indirect communications, such as media coverage of votes, summaries of party platforms, and secondary reports, which combine what political elites say with how journalists contextualize, amplify, or challenge those messages \citep{karlbergPartisanBrandingMedia2002}. Using Liverpool's boycott of The Sun after the 1989 Hillsborough Disaster as a quasi-experiment, \citet{dasilvaMediaInfluenceSpatial2026} found that reduced exposure to ideological media altered both perceived party positions and voting behavior, with the effect on perceived party positions appearing the stronger of the two. Similarly, in the United States, a bundled news treatment changed perceived distances between typical Democratic and Republican voters, demonstrating perceptual updating through indirect communication of party position \citep{levenduskyDoesMediaCoverage2016}.

However, limited research on party brands has engaged with the notion of indirect forms of communications (e.g., media coverage of politicized issues) as the primary source of perceived party brands. Elite speech analysis using roll-call votes and floor speeches captures official positions but misses how those positions are communicated to publics through media \citep{grimmerTextDataPromise2013}. Media coverage, by contrast, captures the discursive environment in which public opinion forms, reflecting both elite cues and journalistic interpretation. Work that does analyze media coverage of political conflict has largely asked how parties of the direct conflict are portrayed rather than how domestic parties are. \citet{pavlichenkoPolarizationMediaPolitical2022}, for instance, examined how British and American outlets constructed Ukraine and Russia as geopolitical actors. That question is distinct from ours, which concerns how the two United States parties were positioned relative to one another within the same coverage. To address this gap, the present study uses media coverage as the primary data source for tracking perceived party brands on Ukraine aid.

\subsection{Ukraine Aid and Foreign Policy Polarization}

Having established why media coverage serves as a critical site for tracking party brand formation, we now turn to foreign policy, where the extent of partisan polarization remains contested.

\subsubsection{Foreign Policy Polarization: Evidence and Debates}

The claim that politics stops at the water's edge has long organized debate over whether foreign policy is insulated from domestic partisan conflict \citep{kupchanDeadCenterDemise2007}. On one hand, for much of the post-World War II era, the Republican and Democratic leadership shared fundamental assumptions about American international engagement, even when disagreeing on specific policies \citep{smeltzAreWeDrowning2022}. On the other hand, a growing body of research documents increasing partisan division on international issues. Studies focusing primarily on congressional behavior and public attitudes find that Democrats and Republicans have grown further apart on foreign policy, mirroring patterns observed in domestic politics \citep{delaetTreatyMakingPartisanPolitics2006,jeongDivisionWatersEdge2019,kupchanDeadCenterDemise2007,schultzPerilsPolarizationUS2017,smeltzAreWeDrowning2022,snyderFreeHandAbroad2009}. Nonetheless, the findings are not consistent. Other scholarship points to continuing areas of agreement or overlap between the parties on foreign policy \citep{busbyComingTogetherComing2020,chaudoinNotOutLiberal2018,scottNotSoSilentPartnerPatterns2014}. \citet{bryanPrevalenceBipartisanshipUS2022} find that extreme partisan division on foreign policy remains the exception and that international issues are still less polarized than domestic issues, at least on congressional floors. Research on both elite and mass attitudes also suggests that the left-right ideological division predicts partisan disagreement more strongly on domestic matters than on international affairs \citep{baldassarriPartisansConstraintPolitical2008a,noelPoliticalIdeologiesPolitical2014}. Foreign policy debates frequently blur partisan cleavages with divisions within parties as well as between them \citep{pratherValuesWatersEdge2014,rathbunWedgesWidgetsLiberalism2016}. Such distinction also changes how apparent consensus should be interpreted. Interparty agreement can coexist with intraparty division, and aggregate party differences can obscure coalitions that cut across party lines without accounting for intraparty agreement and division \citep{friedrichsPolarizationUSForeign2022}.

Additionally, the durability of bipartisan foreign policy consensus carries significant implications beyond academic interest. \citet{bryanPrevalenceBipartisanshipUS2022} argue, partisan division on foreign policy undermines America's capacity to maintain reliable partnerships abroad, honor long-term international obligations, project unified national purpose, fund necessary global initiatives, and adapt strategy based on past experiences. Work on the Ukraine war makes those costs concrete. \citet{stewartRevisitingGlobalPosture2023} assesses what shifting support implies for United States force posture in Europe, and \citet{thakkarRussiaUkraineWarIndia2024} shows that third parties read United States commitment as a signal, examining how India's strategic proximity to Washington interacts with its reluctance to take sides. Understanding when and how foreign policy issues polarize thus addresses questions of both scholarly and practical significance.

\subsubsection{Ukraine Aid as a Particular Case}
United States assistance to Ukraine presents a question about the durability of political alignment. In the first year after Russia's full-scale invasion, policy combined extensive military, economic, diplomatic, and intelligence support with substantial bipartisan backing \citep{harrisWarUkrainePolarised2023a}. Yet early cooperation did not eliminate disagreement. Republican criticism was more visible, a small group called for ending support, and a broader group questioned unrestricted assistance \citep{harrisWarUkrainePolarised2023a}. Congress nevertheless approved large aid package through cross-party coalitions, and later votes continued to pair unified or near-unified Democratic support with substantial, though divided, Republican support \citep{tamaBipartisanshipPutTest2025}. The initial condition was a broad but unsettled support.

By late 2023, Ukraine aid had become a defining partisan issue, featuring prominently in Republican primary debates and contributing to the House Speaker crisis of October 2023. This transformation may reflect broader changes in the conditions enabling foreign policy bipartisanship. Some of those conditions were specific to Russia policy. \citet{podrazaInconsistencyUSPolicy2025} traces how United States policy toward Russia varied across administrations without clear consistency, so the post-invasion consensus rested on agreement about an immediate response rather than on a long-term plan. Additionally, partisan sorting means foreign policy preferences now correlate more strongly with domestic ideology than in previous decades. Media fragmentation enables distinct partisan narratives to develop without common informational foundations. And congressional behavior is increasingly driven by sense to vote with co-partisans regardless of their own policy views \citep{leeInsecureMajoritiesCongress2016,theriaultGingrichSenatorsRoots2013}. The Trump era intensified these dynamics, particularly regarding skepticism toward international commitments.

The Subsequent debate also became more issue-specific. Republican contestation was more pronounced over Ukraine aid than over NATO enlargement or sanctions on Russia \citep{riegerManchurianCandidateGOPs2025}. Cross-case comparison also shows that intraparty division was located differently across aid debates: Republican disagreement centered more heavily on Ukraine, Democratic disagreement on Israel, and Taiwan assistance retained broad backing in both parties \citep{tamaBipartisanshipPutTest2025}. Contemporary Ukraine-aid debate contains distinguishable components rather than one opposition-support continuum. Therefore, to better capture systematically the temporal trajectory of the polarization of Ukraine aid, we must first identify the ideological component in which the issue is polarized on.

\subsubsection{Discursive Formations and Dimensional Structure}

\citet{riegerTheyShouldntReceive2025} identifies four discursive formations that structure contemporary foreign policy contestation: Isolationism, Conservative Realism, Neoconservatism, and Liberal Internationalism. Internationalism emphasizes establishing a stable international order with the United States at its center, including through defense of democratic values against authoritarian threats \citep{coxUSForeignPolicy2018}. Neoconservatism shares internationalist goals but favors unilateral initiatives, viewing commitments like Ukraine support as projections of American power rather than multilateral cooperation \citep{khongNeoconservatismDomesticSources2008}. Isolationism prioritizes American interests through minimizing foreign commitments and their associated costs. Historically, Democrats have generally emphasized multilateral cooperation and strengthening global governance structures, while Republicans have tended to prioritize American autonomy in international decision-making and robust military capabilities as the foundation of national security \citep{smeltzDividedWeStand2020}. \citet{meadJacksonianRevoltAmerican2017} observes, contemporary isolationism often also includes a populist view that frames international commitments as serving elite interests rather than ordinary Americans, redirecting threat perception inward. This skeptical framing enables isolationists to recast issues like Ukraine aid as competing with domestic priorities rather than as standalone foreign policy questions.

Drawing on these formations, we use a deliberately bounded rather than exhaustive set of four dimensions (i.e., a main dimension of general aid stance and three explanatory dimensions). The selection rule for the explanatory dimension is to separate non-overlapping and non-interchangeable political judgments that can vary independently given the question of whether assistance should continue: how the recipient is evaluated, which external threat receives priority, and whether foreign commitments are pitted against domestic priorities. General Ukraine Aid Stance (Q1) captures the general stance from support to opposition. Ukraine Characterization (Q2) isolates portrayals of Ukraine as a democratic partner from portrayals of it as unreliable, burdensome, or outside United States concern; it tracks how the recipient of the aid is evaluated alongside the decision to support/not support aid. Threat Priority (Q3) distinguishes emphasis on Russia and the defense of Ukraine from emphasis on China and the Indo-Pacific \citep{harrisWarUkrainePolarised2023a,riegerManchurianCandidateGOPs2025}. Resource Allocation (Q4) captures whether assistance is represented as competing with domestic welfare, border protection, or other national priorities.

The dimension of Resource Allocation is particularly important in understanding the temporal progression of the general Ukraine aid stance between two parties. \citet{jeongDivisionWatersEdge2019} argue that foreign policy polarization is largely derivative of broader domestic polarization, with ideological conflict on domestic issues spilling over into international affairs. \citet{borgBattleSoulThis2024} extends this reasoning, arguing that foreign policy contestation fundamentally involves competing visions of American identity and priorities. If foreign policy polarization operates primarily through domestic spillover, then framing Ukraine aid as competing with domestic spending directly activates existing partisan divisions over resource allocation, taxation, and government priorities. Notably, this linkage is not exclusive to either party: the Biden administration explicitly connected foreign and domestic policy in its 2021 Interim National Security Strategic Guidance, while Republican critics have framed Ukraine aid as diverting resources from border security and domestic needs. The resource allocation dimension thus provides the most direct bridge between foreign policy and the domestic ideological conflicts that potentially drive the politicization and polarization of Ukraine aid.

\subsubsection{Asymmetric Partisan Incentives}

Furthermore, minority parties often face incentives to challenge incumbent positions on emerging issues, creating space for future contestation. By establishing that an issue is moveable, challengers can claim it as distinctive party territory \citep{abou-chadiModerateNecessaryRole2016,adamsCompanyYouKeep2016,meguidPartyCompetitionUnequals2010}. Over time, this enables challenger parties to establish themselves as the authoritative voice on particular issues and gain outsized electoral returns when those issues become salient \citep{abou-chadiModerateNecessaryRole2016}. \citet{kollbergChallengerAdvantageHow2025} also finds that legislators from challenger parties tend to criticize ``the establishment'' during campaign periods and actively stake out distinctive positions on emerging issues during moments of national uncertainty. Overall, this suggests asymmetric incentives: the minority party benefits from polarizing novel issues to differentiate its brand, while the incumbent party, such as the Democratic Party at the time of 2022, faces costs from abandoning established positions.

While Republicans are not a small or niche party, they faced analogous dynamics as the minority party seeking voter resources to win subsequent elections. Democrats, by contrast, faced potential costs from abandoning a Biden administration signature policy. These asymmetric incentives generate specific predictions: Republicans should drive polarization and party brand updates through positional movement and message consolidation, while Democrats should solidify their existing positioning. Furthermore, if polarization serves party branding functions as theorized, increasing polarization should be accompanied with increasing consistency and coherence of intraparty messaging even as represented by news. As within-party messaging becomes more coherent and between-party positions more distinct, voters relying on partisan heuristics should find it easier to discern what each party stands for on Ukraine aid. We test these predictions empirically, examining not only whether polarization occurred but which party's behavior drove observed polarization and whether within-party coherence increased alongside polarization.

\subsection{Computational Methods for Tracking Party Brand}

Traditional approaches for measuring polarization and perceived coherence of party messages present constraints for tracking discourse over the long term. Survey-based measures such as feeling thermometers capture partisan perception and affect at discrete time points but cannot track continuously without repeated data collection, which incurs substantial financial costs \citep{iyengarOriginsConsequencesAffective2019,pewresearchcenterPoliticalPolarizationAmerican2014}. Survey measures are subject to effects that may inflate polarization estimates, as self-administered online surveys tend to produce higher affective polarization scores than face-to-face interviews \citep{thorntonImpactSurveyMode2025}. While qualitative content analysis provides interpretive depth that identifies nuanced partisan framings and context from political debates and public discourses \citep{krippendorffContentAnalysisIntroduction2019}, traditional human coding struggles with scale. \citet{krippendorffMeasuringReliabilityQualitative2004} notes that inter-coder reliability becomes increasingly difficult to maintain as corpus size grows, and the financial and temporal costs of training multiple coders often constrain sample sizes. For our purpose of tracking the temporal trend of an ongoing issue contestation, both approaches would therefore be costly and still leave the trend under-observed.

Computational text analysis enables systematic examination of large corpora while maintaining measurement consistency across time and sources \citep{grimmerTextDataNew2022}. It leverages existing secondary data, and it permits continuous temporal tracking at granularities that would be expensive with repeated surveys. For instance, \citet{hartPoliticizationPolarizationCOVID192020} measured politicization and polarization in United States newspaper and television coverage of COVID-19 using computational content analysis, establishing that partisan structure in news coverage can be recovered at scale. They covered three months; but our questions concern a trajectory across 34 months. We also have three other requirements: First, we need a position on a dimension we specify in advance rather than a dimension the method discovers, because the dimensions come from the discursive formations reviewed above and not from the corpus. Second, we need that position recovered from a single sentence, because party attribution can be cleanly mapped at sentence level. Third, we need it at a scale of tens of thousands of sentences, and fourth, we need it repeated across 34 months so that the measure is comparable from one month to the next. These four requirements are what this study's text analysis method has to satisfy.

Established methods each satisfy some of them. Wordfish \citep{slapinScalingModelEstimating2008} extracts latent ideological positions from word frequency patterns, identifying the primary dimension of textual variation through unsupervised learning, and scaling of this kind has been used to estimate party positions from election manifestos \citep{backlundSwedenDemocratsPolitical2011}. What it recover depends on assumptions that the output itself does not reveal \citep{egerodScalingPoliticalPositions2020}. Wordfish recovers the dimension along which texts vary most rather than the dimension we specify, and it recovers one dimension rather than several. For Ukraine aid discourse, where polarization occurs potentially along distinct dimensions, that single-dimension condition is limiting. \citet{gentzkowMeasuringGroupDifferences2019a} demonstrate that partisan differentiation in Congress is driven more by changes in how parties frame issues, for example, ``death tax'' versus ``estate tax'', than by changes in which issues they emphasize, suggesting that capturing polarization requires multidimensional measurement of how an issue is characterized rather than simply whether it is discussed. Nonetheless, \citet{koljonenComparingComputationalNoncomputational2022} does reveal that computational and non-computational party position estimates does not necessarily agree, at least in the case of Finish elections from 2003 to 2019. This is a caution that we take up when specifying our own validation. Topic models that accept researcher-supplied keywords address part of the dimension problem, but they still operate on word counts within documents and return the prevalence of a topic rather than a position along it, which is the quantity our research questions require \citep{eshimaKeywordAssistedTopicModels2024}.

Dictionary and embedding approaches do work from predetermined categories, and prior work has applied them to partisan framing in congressional speech \citep{vanmatreNaturalLanguageProcessing2025}. Their performance depends heavily on the corpus and the specification chosen, and \citet{rodriguezWordEmbeddingsWhat2022} show that embedding-based measures require deliberate validation rather than transferring reliably across applications. Both families also are limited at inferring meaning from surrounding context. Supervised classification satisfies the first three requirements and has been used to measure polarization directly: \citet{petersonClassificationAccuracySubstantive2018} treat the accuracy with which a classifier separates the speeches of opposing parties as itself a measure of how far apart those parties stand. That design does not transfer to our case. It requires party labels on the training texts, whereas we deliberately mask party references so that the annotator cannot infer a position from party membership, and it returns a single scalar per period rather than positions on dimensions we specify in advance. training four dimension-specific classifiers would require a labelled corpus large enough to reintroduce the problem of cost that made qualitative analysis infeasible at this scale \citep{ornsteinHowTrainYour2025,wilkersonLargeScaleComputerizedText2017}.

Large language models satisfy all four. They process semantic context, which allows classification of short units, and they can be directed at several predetermined dimensions at once rather than recovering a single latent one. Emerging work establishes that prompted models perform these tasks credibly on political text. \citet{ornsteinHowTrainYour2025} show across sentiment classification, ideology scaling, and topic modeling that prompted models outperform conventional supervised approaches and perform comparably to crowd coders at a fraction of the cost, and \citet{gilardiChatGPTOutperformsCrowd2023} report comparable gains against human annotators. Two further findings speaks directly to our design. \citet{tornbergLargeLanguageModels2025} annotated the political affiliation of politicians from a single social media message across eleven countries and found that an prompt-tuned model exceeded both supervised classifiers and human coders in every language and country context, reaching an accuracy of 0.93 in the United States case. The study further reports that GPT-4, unlike the supervised classifiers, correctly annotated messages requiring interpretation of implicit references or reasoning from contextual knowledge, which is the capacity a single sentence most demands. \citet{mensPositioningPoliticalTexts2025}, separately, show that a model can be prompted directly for a continuous position estimate rather than for a discrete label. Our instrument requires both properties at once, since it annotates a single sentence and returns a position on a scale.

The literature that establishes this capability does not, however, specify how to build the entire LLM annotation suite. Performance varies with prompt wording, supplied context, and model \citep{ornsteinHowTrainYour2025}. General guidance exists on how to construct prompts for text analysis \citep{tornbergBestPracticesText2024}. However, specific validated protocol for evaluating party position on multiple predetermine foreign policy dimension from a single news sentence does not exist. Where the available guidance does promote a choice, it does not always promote a consistent one. \citet{ornsteinHowTrainYour2025} recommend few-shot prompts and find that additional examples improve performance, whereas \citet{tornbergLargeLanguageModels2025} obtains its results zero-shot, from instructions written in natural language and no labelled examples. We follow the latter suggestion, and we set out our reasons alongside the prompt in the Method section.

In the Method section, we address three things. First, we specify our annotation suite, report the model, the request settings, and the exact prompt. Second, we justify each choice on the general principles the literature does supply, since a task-specific standard is unavailable. Third, we validate. \citet{pangakisAutomatedAnnotationGenerative2023} argue that automated annotation with generative models requires validation against human labels on the specific task rather than an assumption of transferable performance, and \citet{ornsteinHowTrainYour2025} treat comparison against human coding as ``golden standard'' against bias introduced during prompt design. Validation against human labels alone, however, would not tell us whether our annotation instructions are reproducible, nor whether the model returns the same score on repeated presentation of the same sentence. We therefore evaluate three properties rather than one: agreement between human coders, which tests whether our instructions are reproducible by people; agreement between the model and human coders, which tests the annotation against human reference; and the model's agreement with itself across repeated trials, which tests reproducibility of the model's own output.

Finally, as \citet{nemethScopingReviewUse2023} cautions, spurious results may arise when causal relations are inferred from textual data. Our research concern the presence, timing and the correlational nature of polarization and the coherence of party elites on a given issue. Given that our data is limited to the supply side, our observational design cannot establish whether the temporal progression observed induces a measured downstream voter effect. Whether the progression of perceived party brand, polarization, and within-party coherence correlates with or leads to changes in partisanship and voting behavior remains a separate empirical matter.

\subsection{Hypothesis}

The preceding review motivates a longitudinal account of perceived party brand that separates temporal development, party contribution, cross-dimensional alignment, coherence of within-party message. We therefore organize the empirical investigation by the following research questions and hypotheses.

RQ1. How did General Ukraine Aid Stance polarization develop from March 2022 through December 2024?

H1a. Elite discourse on Ukraine aid exhibits increasing polarization from March 2022 to December 2024, with partisan divergence on general aid stance (Q1) significantly greater at the study's end than at its beginning.

H1b. This increase in polarization reflects punctuated shifts rather than a pattern of gradual drift over time.

RQ2a. Which party contributed more to the General Ukraine Aid Stance gap?

H2a. Polarization was driven asymmetrically, with Republicans exhibiting greater positional movement toward opposition than Democrats, consistent with minority party incentives to differentiate on emerging issues.

RQ2b. How did Ukraine Characterization, Threat Priority, and Resource Allocation align contemporaneously with General Ukraine Aid Stance polarization?

H2b. Among the explanatory dimensions (Ukraine Characterization, Threat Priority, and Resource Allocation), Resource Allocation framing aligns most strongly temporally to the polarization on General Aid Stance.

RQ3. How is General Ukraine Aid Stance polarization related to Republican and Democratic message coherence?

H3. As polarization moves on the temporal trajectory, both party's within-party coherence also moves with it contemporaneous.

RQ4. How did the Party Brand Index evolve over the study period, and what temporal changes characterized its trajectory?

Overall, our empirical point of interest is how parties are represented in news sentences, as annotated by LLMs and by the four aforementioned dimensions. As we detail later in Method, we build the necessary empirical measurement in three stages. We annotated individual sentences on four dimensions, averaged those annotations within each article so that an article contributes one observation per party and dimension rather than one observation per sentence, and then aggregated articles into monthly series.

